% Options for packages loaded elsewhere
\PassOptionsToPackage{unicode}{hyperref}
\PassOptionsToPackage{hyphens}{url}
\PassOptionsToPackage{space}{xeCJK}
\documentclass[
  a4paper,11pt]{ltjsarticle}
\usepackage{xcolor}
\usepackage[margin=1in]{geometry}
\usepackage{amsmath,amssymb}
\setcounter{secnumdepth}{-\maxdimen} % remove section numbering
\usepackage{iftex}
\ifPDFTeX
  \usepackage[T1]{fontenc}
  \usepackage[utf8]{inputenc}
  \usepackage{textcomp} % provide euro and other symbols
\else % if luatex or xetex
  \usepackage{unicode-math} % this also loads fontspec
  \defaultfontfeatures{Scale=MatchLowercase}
  \defaultfontfeatures[\rmfamily]{Ligatures=TeX,Scale=1}
\fi
\usepackage{lmodern}
\ifPDFTeX\else
  % xetex/luatex font selection
  \ifXeTeX
    \usepackage{xeCJK}
    \setCJKmainfont[]{Hiragino Mincho ProN}
  \fi
  \ifLuaTeX
    \usepackage[]{luatexja-fontspec}
    \setmainjfont[]{Hiragino Mincho ProN}
  \fi
\fi
% Use upquote if available, for straight quotes in verbatim environments
\IfFileExists{upquote.sty}{\usepackage{upquote}}{}
\IfFileExists{microtype.sty}{% use microtype if available
  \usepackage[]{microtype}
  \UseMicrotypeSet[protrusion]{basicmath} % disable protrusion for tt fonts
}{}
\makeatletter
\@ifundefined{KOMAClassName}{% if non-KOMA class
  \IfFileExists{parskip.sty}{%
    \usepackage{parskip}
  }{% else
    \setlength{\parindent}{0pt}
    \setlength{\parskip}{6pt plus 2pt minus 1pt}}
}{% if KOMA class
  \KOMAoptions{parskip=half}}
\makeatother
\usepackage{longtable,booktabs,array}
\newcounter{none} % for unnumbered tables
\usepackage{calc} % for calculating minipage widths
% Correct order of tables after \paragraph or \subparagraph
\usepackage{etoolbox}
\makeatletter
\patchcmd\longtable{\par}{\if@noskipsec\mbox{}\fi\par}{}{}
\makeatother
% Allow footnotes in longtable head/foot
\IfFileExists{footnotehyper.sty}{\usepackage{footnotehyper}}{\usepackage{footnote}}
\makesavenoteenv{longtable}
\setlength{\emergencystretch}{3em} % prevent overfull lines
\providecommand{\tightlist}{%
  \setlength{\itemsep}{0pt}\setlength{\parskip}{0pt}}
\usepackage{bookmark}
\IfFileExists{xurl.sty}{\usepackage{xurl}}{} % add URL line breaks if available
\urlstyle{same}
\hypersetup{
  hidelinks,
  pdfcreator={LaTeX via pandoc}}

\author{}
\date{}

\begin{document}

\section{世代混合の幾何学的構造------CKM行列とPMNS行列の定性的導出}\label{ux4e16ux4ee3ux6df7ux5408ux306eux5e7eux4f55ux5b66ux7684ux69cbux9020ckmux884cux5217ux3068pmnsux884cux5217ux306eux5b9aux6027ux7684ux5c0eux51fa}

\subsection{------ {[}4{]} §9.11.6 の世代混合問題に対する構造的回答
------}\label{ux306eux4e16ux4ee3ux6df7ux5408ux554fux984cux306bux5bfeux3059ux308bux69cbux9020ux7684ux56deux7b54}

\textbf{著者}: 木原 範昭 (WF System Co., Ltd.)

\textbf{日付}: 2026年4月

\begin{center}\rule{0.5\linewidth}{0.5pt}\end{center}

\subsection{§0 本稿の目的}\label{ux672cux7a3fux306eux76eeux7684}

{[}4{]} §9.11.6 において、「世代間の混合（CKM 行列、PMNS
行列）が、\(xyz\) 軸の \(S_3\)
対称性の破れとしてどのように導出されるかは未解決である」と宣言した。

本稿はこの宣言に応じ、以下を示す。

\begin{enumerate}
\def\labelenumi{\arabic{enumi}.}
\tightlist
\item
  CKM行列の対角優位構造が、ヒッグス（定常波）の軸選択による
  \(\mathrm{SO}(3) \to \mathrm{SO}(2)\)
  対称性の破れから直接従うこと（§1）
\item
  PMNS行列の大混合角構造が、ニュートリノの（空間,
  空間）型2-面構造から直接従うこと（§2）
\item
  CKMとPMNSの構造的差異が、（空間, \(t\)）型と（空間,
  空間）型の幾何学的差異に帰着すること（§3）
\end{enumerate}

本稿の主張は混合行列の\textbf{定性的構造}（対角優位 vs
大混合角）の導出に限定される。混合角の具体的数値の導出は、質量比の定量的導出（{[}4{]}
§9.11.4）と同根の未解決問題であり、本稿の範囲外である。

\textbf{参照論文}: - {[}4{]}
6次元超直方体の集合構造とその組合せ論的性質（W8） - {[}5{]}
形不変波の相互作用（W11）

\begin{center}\rule{0.5\linewidth}{0.5pt}\end{center}

\subsection{\texorpdfstring{§1 CKM行列: （空間,
\(t\)）型の世代混合}{§1 CKM行列: （空間, t）型の世代混合}}\label{ckmux884cux5217-ux7a7aux9593-tux578bux306eux4e16ux4ee3ux6df7ux5408}

\subsubsection{1.1
世代と空間軸の対応}\label{ux4e16ux4ee3ux3068ux7a7aux9593ux8ef8ux306eux5bfeux5fdc}

{[}4{]} §9.1
により、荷電フェルミオン（クォークおよび荷電レプトン）は（空間1軸,
\(t\)）型の2-面を持つ。3つの空間軸 \(x, y, z\) が3世代に対応する:

{\def\LTcaptype{none} % do not increment counter
\begin{longtable}[]{@{}cccc@{}}
\toprule\noalign{}
世代 & 空間軸 & 上型クォーク & 下型クォーク \\
\midrule\noalign{}
\endhead
\bottomrule\noalign{}
\endlastfoot
1 & \(x\) & \(u = (+, 0, 0, +, 0, Q)\) & \(d = (+, 0, 0, -, 0, Q)\) \\
2 & \(y\) & \(c = (0, +, 0, +, 0, Q)\) & \(s = (0, +, 0, -, 0, Q)\) \\
3 & \(z\) & \(t = (0, 0, +, +, 0, Q)\) & \(b = (0, 0, +, -, 0, Q)\) \\
\end{longtable}
}

\subsubsection{1.2
ヒッグスの軸選択と対称性の破れ}\label{ux30d2ux30c3ux30b0ux30b9ux306eux8ef8ux9078ux629eux3068ux5bfeux79f0ux6027ux306eux7834ux308c}

{[}4{]} §9.11.7 により、ヒッグスボソン（定常波、\(\kappa = 1\)）は
\(xyz\) 3軸のうち1軸に非零成分を取る。初期状態では \(xyz\) は
\(\mathrm{SO}(3)\)
対称性のもとで完全に対等であるが、ヒッグスが1軸を選ぶことで対称性が自発的に破れる。選ばれた軸を
\(z\) と表記する。

この結果: - \(z\) 軸: ヒッグスの非零軸 → 直接結合 → 質量最大（第3世代）
- \(x, y\) 軸: ヒッグスの非零軸を共有しない → 間接結合のみ →
質量が軽い（第1・第2世代）

対称性は \(\mathrm{SO}(3) \to \mathrm{SO}(2)\) に破れ、\(z\)
軸が分離する一方、\(x\) 軸と \(y\) 軸の間には近似的な \(\mathrm{SO}(2)\)
対称性が残存する。

\subsubsection{1.3
CKM行列の対角優位構造}\label{ckmux884cux5217ux306eux5bfeux89d2ux512aux4f4dux69cbux9020}

\textbf{命題 1.1}. （空間,
\(t\)）型フェルミオンの世代混合において、ヒッグスの \(z\)
軸選択は以下の階層構造を導く:

\begin{enumerate}
\def\labelenumi{(\roman{enumi})}
\item
  第3世代と第1・第2世代の混合は小さい（\(z\) 軸の分離）
\item
  第1世代と第2世代の混合は (i) より大きい（\(x, y\) 軸の近似的対称性）
\item
  CKM行列は対角優位である
\end{enumerate}

\emph{論証}. \(W^{\pm}\) による弱相互作用はアイソスピン遷移（{[}5{]}
§6）であり、軸型対称性により位置型軸全体に作用してフェルミオンの \(t\)
軸符号を反転させ、上型クォークを下型クォークに変換する。各世代のクォークは1本の空間軸のみを占有するため、世代間遷移は空間軸の変更を伴う。

ヒッグスの \(z\) 軸選択により、\(z\)
軸上のフェルミオン（第3世代）は他の軸上のフェルミオンと質量スケールが大きく異なる。質量固有状態とフレーバー固有状態のずれ（混合角）は、軸間の結合強度に依存し、結合強度は質量差に反比例する。

\begin{enumerate}
\def\labelenumi{(\roman{enumi})}
\item
  \(m_3 \gg m_{1,2}\) より、第3世代と第1・第2世代の間の混合角は小さい。
\item
  \(x\) 軸と \(y\) 軸は \(\mathrm{SO}(2)\)
  対称性のもとで近似的に対等であるため、\(m_1 \approx m_2\)
  であり、第1・第2世代間の混合角は (i) より大きい。
\item
  (i), (ii)
  より、CKM行列は対角成分が1に近く、非対角成分が小さい（対角優位）。\(\square\)
\end{enumerate}

\textbf{実験値との比較}:

\[V_{\text{CKM}} \approx \begin{pmatrix} 0.974 & 0.225 & 0.004 \\ 0.225 & 0.973 & 0.041 \\ 0.009 & 0.040 & 0.999 \end{pmatrix}\]

\begin{itemize}
\tightlist
\item
  \(|V_{ub}| = 0.004\), \(|V_{td}| = 0.009\): 第3↔第1世代の混合が最小 ✓
\item
  \(|V_{cb}| = 0.041\), \(|V_{ts}| = 0.040\): 第3↔第2世代の混合が小 ✓
\item
  \(|V_{us}| = 0.225\): 第1↔第2世代の混合が最大 ✓
\end{itemize}

\begin{center}\rule{0.5\linewidth}{0.5pt}\end{center}

\subsection{§2 PMNS行列: （空間,
空間）型の世代混合}\label{pmnsux884cux5217-ux7a7aux9593-ux7a7aux9593ux578bux306eux4e16ux4ee3ux6df7ux5408}

\subsubsection{2.1
ニュートリノの2-面構造}\label{ux30cbux30e5ux30fcux30c8ux30eaux30ceux306e2-ux9762ux69cbux9020}

{[}4{]} §9.9 表Aにより、ニュートリノは（空間,
空間）型の2-面を持つ。\(\binom{3}{2} = 3\)
通りの空間軸の組合せが3フレーバーに対応する:

{\def\LTcaptype{none} % do not increment counter
\begin{longtable}[]{@{}ccl@{}}
\toprule\noalign{}
フレーバー & 非零軸 & 集合 \\
\midrule\noalign{}
\endhead
\bottomrule\noalign{}
\endlastfoot
\(\nu_e\) & \((x, y)\) & \((+, +, 0, 0, 0, 0)\) \\
\(\nu_\mu\) & \((y, z)\) & \((0, +, +, 0, 0, 0)\) \\
\(\nu_\tau\) & \((x, z)\) & \((+, 0, +, 0, 0, 0)\) \\
\end{longtable}
}

\subsubsection{2.2 軸共有構造}\label{ux8ef8ux5171ux6709ux69cbux9020}

\textbf{命題 2.1}.
任意の2つのニュートリノフレーバーは、ちょうど1本の空間軸を共有する。

\emph{証明}. 3元集合 \(\{x, y, z\}\) から2元を選ぶ3つの組合せ
\(\{x,y\}, \{y,z\}, \{x,z\}\)
において、任意の2組合せの共通部分はちょうど1元である:

{\def\LTcaptype{none} % do not increment counter
\begin{longtable}[]{@{}ccc@{}}
\toprule\noalign{}
ペア & 共有軸 & 非共有軸 \\
\midrule\noalign{}
\endhead
\bottomrule\noalign{}
\endlastfoot
\(\nu_e, \nu_\mu\) & \(y\) & \(x \leftrightarrow z\) \\
\(\nu_\mu, \nu_\tau\) & \(z\) & \(y \leftrightarrow x\) \\
\(\nu_e, \nu_\tau\) & \(x\) & \(y \leftrightarrow z\) \\
\end{longtable}
}

これは \(\binom{3}{2}\)
の組合せ論的性質であり、3軸・2選択の場合に一般に成立する。\(\square\)

\subsubsection{2.3
PMNS行列の大混合角構造}\label{pmnsux884cux5217ux306eux5927ux6df7ux5408ux89d2ux69cbux9020}

\textbf{命題 2.2}. （空間,
空間）型フェルミオンの世代混合において、ヒッグスの \(z\)
軸選択は以下の構造を導く:

\begin{enumerate}
\def\labelenumi{(\roman{enumi})}
\item
  \(\nu_\mu\) と \(\nu_\tau\) はともに \(z\)
  軸（ヒッグス軸）を含み、非共有軸 \((y, x)\)
  はヒッグスの影響を受けないため、\(\nu_\mu \leftrightarrow \nu_\tau\)
  混合はほぼ最大となる
\item
  \(\nu_e\) は \(z\) 軸を含まないため、\(\nu_e \leftrightarrow \nu_\mu\)
  および \(\nu_e \leftrightarrow \nu_\tau\) の混合は (i) より小さい
\item
  PMNS行列は非対角成分が大きい（CKMと対照的）
\end{enumerate}

\emph{論証}.

\begin{enumerate}
\def\labelenumi{(\roman{enumi})}
\item
  \(\nu_\mu = (y, z)\) と \(\nu_\tau = (x, z)\) は共有軸 \(z\)
  を持ち、非共有軸はそれぞれ \(y\) と \(x\) である。\(x\) と \(y\)
  はヒッグスの影響を受けず \(\mathrm{SO}(2)\)
  対称性が残存するため、\(\nu_\mu\) と \(\nu_\tau\)
  のヒッグスとの結合強度はほぼ等しい。結合強度が等しい2状態の混合角は最大（\(\theta \approx 45°\)）に近づく。
\item
  \(\nu_e = (x, y)\) は \(z\) 軸を含まない。\(\nu_e\) とヒッグスの結合は
  \(x, y\) 軸を経由する間接的なものに限られ、\(z\) 軸を含む
  \(\nu_\mu, \nu_\tau\) とは結合構造が異なる。この非対称性により
  \(\nu_e\) と \(\nu_{\mu,\tau}\) の混合角は最大値から外れる。
\item
  \begin{enumerate}
  \def\labelenumii{(\roman{enumii})}
  \tightlist
  \item
    の最大混合と (ii)
    の中程度の混合が共存するため、PMNS行列全体として非対角成分が大きい。\(\square\)
  \end{enumerate}
\end{enumerate}

\textbf{実験値との比較}:

\begin{itemize}
\tightlist
\item
  \(\theta_{23} \approx 49°\)（大気角、\(\nu_\mu \leftrightarrow \nu_\tau\)）:
  ほぼ最大 ✓ --- 命題 (i)
\item
  \(\theta_{12} \approx 33°\)（太陽角、\(\nu_e \leftrightarrow \nu_\mu\)）:
  大きいが最大ではない ✓ --- 命題 (ii)
\item
  \(\theta_{13} \approx 8.6°\)（原子炉角、\(\nu_e \leftrightarrow \nu_\tau\)）:
  最小 ✓ --- 命題 (ii)
\end{itemize}

\begin{center}\rule{0.5\linewidth}{0.5pt}\end{center}

\subsection{§3
構造的差異の起源}\label{ux69cbux9020ux7684ux5deeux7570ux306eux8d77ux6e90}

\subsubsection{\texorpdfstring{3.1 （空間, \(t\)）型 vs （空間,
空間）型}{3.1 （空間, t）型 vs （空間, 空間）型}}\label{ux7a7aux9593-tux578b-vs-ux7a7aux9593-ux7a7aux9593ux578b}

\textbf{命題 3.1}.
CKM行列（対角優位）とPMNS行列（大混合角）の構造的差異は、2-面の軸構成の違いに帰着する。

{\def\LTcaptype{none} % do not increment counter
\begin{longtable}[]{@{}
  >{\raggedright\arraybackslash}p{(\linewidth - 4\tabcolsep) * \real{0.3333}}
  >{\raggedright\arraybackslash}p{(\linewidth - 4\tabcolsep) * \real{0.3333}}
  >{\raggedright\arraybackslash}p{(\linewidth - 4\tabcolsep) * \real{0.3333}}@{}}
\toprule\noalign{}
\begin{minipage}[b]{\linewidth}\raggedright
性質
\end{minipage} & \begin{minipage}[b]{\linewidth}\raggedright
（空間, \(t\)）型
\end{minipage} & \begin{minipage}[b]{\linewidth}\raggedright
（空間, 空間）型
\end{minipage} \\
\midrule\noalign{}
\endhead
\bottomrule\noalign{}
\endlastfoot
占有する空間軸 & \textbf{1本} & \textbf{2本} \\
ヒッグス軸(\(z\))を含む世代 & 第3世代のみ & \(\nu_\mu\) と \(\nu_\tau\)
の\textbf{2つ} \\
2フレーバー間の軸共有 & なし（各世代が独立軸） &
常に1軸を共有（命題2.1） \\
ヒッグスによる分離 & 1 vs 2 に分離（\(z\) vs \(x,y\)） &
弱い分離（\(z\)を含む2つ vs 含まない1つ） \\
混合角の帰結 & 小さい（対角優位） & 大きい \\
\end{longtable}
}

\emph{論証}. （空間,
\(t\)）型では各世代が1本の独立な空間軸を占有し、世代間に軸の共有がない。ヒッグスの
\(z\)
軸選択は1世代のみを分離し、残る2世代間の結合も間接的であるため、混合角は全体的に小さい。

（空間,
空間）型では各フレーバーが2本の空間軸を占有し、任意の2フレーバーが1軸を共有する（命題2.1）。ヒッグスの
\(z\)
軸は3フレーバー中2つに含まれるため分離効果が弱く、軸共有による直接的結合が混合角を大きくする。\(\square\)

\subsubsection{3.2
補講1との関係}\label{ux88dcux8b1b1ux3068ux306eux95a2ux4fc2}

補講1 §2.2 で示したように、（空間, 空間）型のニュートリノは
\(\text{sign}(M_F) = +1\)（正配向）であり、（空間,
\(t\)）型の荷電フェルミオンとは配向が異なる。補講1 命題2.2
により、反ニュートリノも \(M_F^{\text{anti}} = M_F\)（符号不変）を持つ。

世代混合の構造的差異は、この配向の違いと同じ幾何学的起源------（空間,
\(t\)）型と（空間, 空間）型の2-面構造の違い------を持つ。

\begin{center}\rule{0.5\linewidth}{0.5pt}\end{center}

\subsection{§4 本稿の結論}\label{ux672cux7a3fux306eux7d50ux8ad6}

本稿では、{[}4{]} §9.11.6
で未解決とされた世代混合の問題に対し、以下の構造的回答を与えた。

\begin{enumerate}
\def\labelenumi{\arabic{enumi}.}
\item
  \textbf{CKM行列の対角優位構造}は、（空間,
  \(t\)）型フェルミオンの各世代が独立な1空間軸を占有し、ヒッグスの \(z\)
  軸選択により \(\mathrm{SO}(3) \to \mathrm{SO}(2)\)
  の対称性の破れが第3世代を分離することから従う（命題1.1）。
\item
  \textbf{PMNS行列の大混合角構造}は、（空間,
  空間）型ニュートリノの各フレーバーが2空間軸を占有し、任意の2フレーバーが1軸を共有する（命題2.1）ことから従う。特に
  \(\nu_\mu\) と \(\nu_\tau\) がともにヒッグス軸 \(z\)
  を含むため、\(\theta_{23}\) がほぼ最大となる（命題2.2）。
\item
  \textbf{CKMとPMNSの構造的差異}は、（空間, \(t\)）型と（空間,
  空間）型の2-面構造の違いに帰着する（命題3.1）。
\end{enumerate}

混合角の具体的数値の導出は、{[}4{]} §9.11.4
の質量比の定量的導出と同根の未解決問題として残る。

\begin{center}\rule{0.5\linewidth}{0.5pt}\end{center}

\subsection{参考文献}\label{ux53c2ux8003ux6587ux732e}

\begin{itemize}
\tightlist
\item
  {[}4{]} 木原 範昭, ``6次元超直方体の集合構造とその組合せ論的性質,''
  Zenodo preprint, DOI: 10.5281/zenodo.19762134 (2026).
\item
  {[}5{]} 木原 範昭,
  ``形不変波の相互作用------軸方向変位転送・波束変形・因果の遡及的構成,''
  Zenodo preprint, DOI: 10.5281/zenodo.19763463 (2026).
\end{itemize}

\end{document}
